% ID: SEL-D0016-BA-E01
% ESTADO: cotejada. Segunda pasada visual completa y borrador compilado revisados.
% Fuente: Selectividad/matematicas_aplicadas/matematicas_aplicadas_P_D.pdf
% SHA-256: 11865b83e509d9cc2723c752f7a81f714378e1b651f4a4d010dab4a4576b74c2
% Página: 1; bloque A; ejercicio original 1.

\textbf{(2.5 puntos)} Un cocinero tiene que hacer el postre para una cena
y le han encargado dos de sus mejores creaciones: \textit{Delicia Roja} y
\textit{Delicia Negra}. Para elaborar 1 kg de \textit{Delicia Roja} son
necesarias 3 tarrinas de fresas y 1 tableta de chocolate y para elaborar 1 kg
de \textit{Delicia Negra} se necesita 1 tarrina de fresas y 2 tabletas de
chocolate. Dispone de 15 tarrinas de fresas y 10 tabletas de chocolate.
Además, la cantidad de \textit{Delicia Negra} no debe ser inferior a 1.5 kg
y tampoco debe ser superior al doble de \textit{Delicia Roja}. Si cada
kilogramo de \textit{Delicia Roja} le reporta un beneficio de 3 euros y el de
\textit{Delicia Negra} 5 euros, averigüe qué cantidad de cada postre debe
elaborar para conseguir un beneficio máximo y a cuánto asciende ese beneficio.
